% ============================================================
%  LATEX PATCH for "Fungal Hyphae as Distributed Hygroscopic Vapor Sinks"
%  Generated: 2026-04-18
%
%  This file contains drop-in replacement blocks keyed to line numbers
%  in the current manuscript .tex source.  Each block is self-contained
%  and can be copy-pasted into the manuscript.
%
%  Numbers below are from the automated tile-based 2D FFT pipeline
%  (spectral_slope_2d.py) with focus-gate ROI filtering.
% ============================================================


% ============================================================
% BLOCK A — RESULTS (replaces lines ~162–169 + INSERT SEM SECTION)
%
% Context: follows the sentence ending "...separates from
%   \textit{Mucor} ($140 \pm 33$~\textmu m) and \textit{Rhizopus}
%   ($110 \pm 23$~\textmu m)."
%
% Replaces everything from "This correspondence suggests..." through
% "(INSERT SEM SECTION)" with three paragraphs.
% ============================================================

This correspondence suggests that colony-scale architecture modulates
how far the depletion zone extends.  To quantify this architectural
difference, we performed automated two-dimensional Fourier analysis of
colony surface micrographs (Fig.~\ref{fig:fungi}G,H).  Each grayscale
image was divided into overlapping $512 \times 512$-pixel tiles
(stride 256~px); out-of-focus tiles (Laplacian variance below the
10th within-image percentile) and textureless tiles (intensity
standard deviation $< 5$ grey levels) were excluded.  For each
retained tile, a 2D Hann window was applied, the two-dimensional power
spectrum was computed, and azimuthally averaged into radial frequency
bins.  The log--log spectral slope~$\alpha$ was obtained per image by
fitting $\log_{10}(\mathrm{power})$ versus
$\log_{10}(\mathrm{frequency})$ on the mean radial profile over
0.01--0.45~cyc/px.  \textit{Aspergillus} colonies exhibited a
shallower spectral rolloff ($\alpha = -2.96 \pm 0.03$, $n = 14$
images) than \textit{Mucor} ($\alpha = -3.16 \pm 0.11$, $n = 7$
images; Welch's $t = 4.6$, $p = 3.1 \times 10^{-3}$; Mann--Whitney
$p = 1.7 \times 10^{-5}$; Cohen's $d = 2.96$; Kolmogorov--Smirnov
$D = 1.0$; Fig.~\ref{fig:fungi}I), indicating greater fine-scale
surface complexity consistent with densely packed conidial heads
versus the smooth filaments of \textit{Mucor} aerial hyphae.  A
sensitivity analysis confirmed that $\alpha$ was invariant to the
stringency of the focus gate (Supplementary Table~X).

Light microscopy of disaggregated colony tissue confirms that the two
genera differ in unit morphology (Fig.~\ref{fig:fungi}G,H).
\textit{Aspergillus} colonies break apart into dense conidial
aggregates that fill much of the field of view as a single dominant
cluster, whereas \textit{Mucor} tissue separates into a distributed
filamentous network of many small, loosely connected fragments.
Quantitative image analysis across three magnifications ($10\times$,
$20\times$, $40\times$; six \textit{Aspergillus} and three
\textit{Mucor} fields) shows that the spatial heterogeneity of
fungal-material distribution, measured as the coefficient of variation
of local patch density, is consistently higher for
\textit{Aspergillus} than for \textit{Mucor} ($0.69 \pm 0.11$ vs.\
$0.39 \pm 0.03$; Mann--Whitney $U = 18$, $p = 0.024$), with the
ratio preserved at every magnification ($1.7\times$ at $10\times$,
$1.7\times$ at $20\times$, $1.9\times$ at $40\times$).  Because these
structural contrasts are preserved across a fourfold range in
magnification, the measured patchiness reflects intrinsic morphology
rather than preparation-dependent material retention.  The ball-like
conidial aggregates of \textit{Aspergillus} introduce high
spatial-frequency content in any imaging plane, producing the
shallower spectral rolloff observed in colony-surface micrographs
(Fig.~\ref{fig:fungi}I), while the smooth, uniformly distributed
hyphae of \textit{Mucor} yield a steeper spectral decay.  This
structural basis explains why colony-scale architecture, rather than
cell-wall chemistry, is the primary determinant of depletion-zone
width.


% ============================================================
% BLOCK B — METHODS: replace "Spectral analysis of colony surface
% morphology" (lines ~333–342)
% ============================================================

\paragraph{Spectral analysis of colony surface morphology.}
Colony surface micrographs were acquired using the same DSLR macro
imaging system at matched magnification (calibrated at ${\sim}1$~\textmu
m/px using a stage micrometer with 200~\textmu m divisions).
Grayscale images of \textit{Aspergillus} (14~images) and
\textit{Mucor} (7~images) colonies were analyzed using an automated
tile-based 2D FFT pipeline.  Each image was divided into overlapping
$512 \times 512$-pixel tiles (stride 256~px).  Tile-level quality
control proceeded by two gates: (i)~tiles below the 10th percentile
of per-image Laplacian variance were rejected as out-of-focus; and
(ii)~tiles with intensity standard deviation below 5~grey levels were
rejected as textureless.  Typical retention was 75--90\% of tiles per
image.  For each retained tile, a 2D Hann window was applied, the
discrete 2D Fourier transform was computed, and the magnitude-squared
power spectrum was azimuthally averaged into radial frequency bins.  A
mean radial power profile was assembled per image by averaging across
all retained tiles (typically 360--470 per image).  The log--log
spectral slope~$\alpha$ was obtained by ordinary least-squares
regression of $\log_{10}(\mathrm{power})$ versus
$\log_{10}(\mathrm{frequency})$ over the range 0.01--0.45~cyc/px.
Because $\alpha$ is the exponent of a power-law fit in frequency
space, it is independent of absolute pixel calibration.  A sensitivity
analysis varying the focus-gate cutoff from the 0th to the 30th
percentile confirmed that $\alpha$ estimates were invariant to gate
stringency (Supplementary Table~X).  Statistical comparison used
Welch's $t$-test (unequal variances confirmed by Levene's test,
$p < 0.01$).  Both slope distributions passed the Shapiro--Wilk
normality test.


% ============================================================
% BLOCK C — METHODS: new subsection (insert after Block B)
% ============================================================

\paragraph{Light microscopy of colony structure.}
Colony tissue was excised from mature cultures, disaggregated on glass
microscope slides, and imaged under brightfield illumination at
$10\times$, $20\times$, and $40\times$ magnification.  Pixel size was
calibrated from a stage micrometer with 10~\textmu m divisions
(0.75~\textmu m\,px$^{-1}$ at $10\times$;
0.38~\textmu m\,px$^{-1}$ at $20\times$;
0.19~\textmu m\,px$^{-1}$ at $40\times$).  Grayscale images were
segmented by Otsu thresholding of the inverted intensity, followed by
morphological opening and closing to remove speckle noise.  Three
structural metrics were extracted from each field: (i)~cluster
dominance, defined as the area of the largest connected component
divided by the total foreground area; (ii)~median pore diameter,
computed from the Euclidean distance transform of the binary mask; and
(iii)~local density coefficient of variation, computed as the standard
deviation divided by the mean of foreground-area fractions measured
across a regular grid of $8 \times 8$ non-overlapping patches.  Six
\textit{Aspergillus} (two per magnification) and three \textit{Mucor}
(one per magnification) fields were analyzed.


% ============================================================
% BLOCK D — DISCUSSION: update spectral numbers
% (replaces the sentence at ~line 236 starting "Spectral analysis
% supports this interpretation:")
% ============================================================

Spectral analysis supports this interpretation:
\textit{Aspergillus} exhibits significantly greater fine-scale surface
roughness ($\alpha = -2.96$ vs.\ $-3.16$; $p = 3.1 \times 10^{-3}$;
Cohen's $d = 2.96$), consistent with densely packed conidial heads
presenting more absorbing surface per unit footprint and sustaining
greater total vapor flux \cite{sun2017}.


% ============================================================
% BLOCK E — STATISTICAL ANALYSIS: update FFT description
% (replaces "110 FFT transects (70 Aspergillus, 40 Mucor)" at ~line 365)
% ============================================================

% OLD: 110 FFT transects (70 \textit{Aspergillus}, 40 \textit{Mucor}).
% NEW:
14 \textit{Aspergillus} and 7 \textit{Mucor} colony-surface images
(automated tile-based 2D FFT; typically 360--470 tiles per image after
quality-control gating).


% ============================================================
% BLOCK F — FIGURE 3 CAPTION: update panels G–J
% (replaces lines starting "\textbf{(G,\,H)}" through "\textbf{(J)}")
% ============================================================

\textbf{(G,\,H)}~Light micrographs of disaggregated colony tissue at
$40\times$ magnification.  \textit{Aspergillus}~(G) shows dense
conidial aggregates; \textit{Mucor}~(H) shows a distributed
filamentous network.  Scale bars: 50~\textmu m.
\textbf{(I)}~Spectral slopes~$\alpha$ from automated tile-based 2D
FFT analysis of colony surface micrographs (14
\textit{Aspergillus}, 7 \textit{Mucor} images).
\textbf{(J)}~$\alpha$ versus $\bar{\delta}$ (group means~$\pm$~1~SD).
\end{filecontents}
